\section{本章小结}
\label{sec:mrel-summary}

本章开头问了两个问题：两个流形何时在光滑意义下相同，一个对象何时是另一个流形中的光滑子集。第一个问题由微分同胚回答。它是光滑且逆光滑的双射，比同胚多保留了光滑结构，比集合双射多保留了拓扑；$x\mapsto x^3$是光滑的同胚，却不是微分同胚。

微分同胚在每点的微分都是线性同构，所以维数相同。反过来，反函数定理只说明微分在一点可逆时映射在该点附近是微分同胚。$E(t)=(\cos t,\sin t)$与圆周的倍角映射都是局部微分同胚，却不是单射，整体的双射性必须另外核对。微分同胚也不保持长度，$(0,1)$与$\mathbb R$微分同胚。

映射的局部结构由微分的秩决定。浸入的微分处处单射，淹没的微分处处满射，它们分别要求源的维数不超过或不小于目标的维数，并且与映射本身是否单射、满射无关。一点的秩只给出附近秩的下限，$x\mapsto x^2$在原点秩为零、附近秩为一，没有局部标准形式。

在秩保持不变的邻域中，常秩定理用反函数定理把映射化为“前$r$个坐标不变、其余坐标为零”的标准形式。满秩是开条件，所以浸入与淹没自动满足常秩条件，分别局部像坐标子空间的包含与坐标投影。

第二个问题需要整体条件。“8”字形曲线是单射浸入，参数却可以逃向端点而像回到原点，它到像的逆不连续；环面上的无理斜率直线则处处如此。嵌入在浸入之上要求到像的同胚，紧流形上的单射浸入自动是嵌入。

嵌入的像正是嵌入子流形，即每点附近有切片坐标、子流形对应一部分坐标为零的集合；切空间是前几个坐标基向量张成的子空间。圆周、球面、函数图像与乘积都有切片坐标，十字形和“8”字的交叉点则没有，尖点虽然拓扑上是直线，也没有切片坐标。

约束方程提供了构造子流形的一般办法。若$c$是$\Phi:M\to N$的正则值，水平集$\Phi^{-1}(c)$是$m-n$维嵌入子流形，切空间是$\ker\mathrm d\Phi_p$。球面与椭球面由此得到；正交群$\mathrm O(n)$也由此得到，但约束映射要取值于对称矩阵空间，以免把重复的方程计入满秩要求，所得维数为$n(n-1)/2$，单位元处的切空间是反对称矩阵。正则值只是充分条件，$(\lVert\bm y\rVert_2^2-1)^2$的零水平集是圆，却没有正则点。

本章始终在无边界情形下讨论，带边界时反函数定理与正则水平集定理都要另核条件。所有结论都只涉及光滑结构，不涉及长度和角度。在第~\ref{chap:riemannian-geometry}章给切空间配上内积之后，才能问一个微分同胚是否保持长度，即两个黎曼流形是否等距。
